\documentclass[12pt, titlepage]{article}

\usepackage{booktabs}
\usepackage{tabularx}
\usepackage{hyperref}
\hypersetup{
    colorlinks,
    citecolor=blue,
    filecolor=black,
    linkcolor=red,
    urlcolor=blue
}
\usepackage[round]{natbib}

\input{../Comments.text}
\input{../Common.text}

\begin{document}

\title{System Verification and Validation Plan for \progname{}} 
\author{\authname}
\date{\today}
	
\maketitle

\pagenumbering{roman}

\section*{Revision History}

\begin{tabularx}{\textwidth}{p{3cm}p{2cm}X}
\toprule {\bf Date} & {\bf Version} & {\bf Notes}\\
\midrule
Date 1 & 1.0 & Notes\\
Date 2 & 1.1 & Notes\\
\bottomrule
\end{tabularx}

~\\
\wss{The intention of the VnV plan is to increase confidence in the software.
However, this does not mean listing every verification and validation technique
that has ever been devised.  The VnV plan should also be a \textbf{feasible}
plan. Execution of the plan should be possible with the time and team available.
If the full plan cannot be completed during the time available, it can either be
modified to ``fake it'', or a better solution is to add a section describing
what work has been completed and what work is still planned for the future.}

\wss{The VnV plan is typically started after the requirements stage, but before
the design stage.  This means that the sections related to unit testing cannot
initially be completed.  The sections will be filled in after the design stage
is complete.  the final version of the VnV plan should have all sections filled
in.}

\newpage

\tableofcontents

\listoftables
\wss{Remove this section if it isn't needed}

\listoffigures
\wss{Remove this section if it isn't needed}

\newpage

\section{Symbols, Abbreviations, and Acronyms}

\renewcommand{\arraystretch}{1.2}
\begin{tabular}{l l} 
  \toprule		
  \textbf{symbol} & \textbf{description}\\
  \midrule 
  T & Test\\
  \bottomrule
\end{tabular}\\

\wss{symbols, abbreviations, or acronyms --- you can simply reference the SRS
  \citep{SRS} tables, if appropriate}

\wss{Remove this section if it isn't needed}

\newpage

\pagenumbering{arabic}

This document ... \wss{provide an introductory blurb and roadmap of the
  Verification and Validation plan}

\section{General Information}

\subsection{Summary}

FPVR is a virtual reality (VR)-based first-person-view (FPV) drone training simulator designed to help beginners develop and improve their drone piloting skills in a safe and accessible environment. Instead of requiring a physical drone, users can operate a simulated drone using VR controls and practise essential flight manoeuvres without the risks of equipment damage or unsafe flying conditions. It is also meant as a tool for experienced Drone fliers to be able to practice in a safe and controlled environment without the pressure of piloting a very expensive piece of equipment.

The system features free-flight practice, progressively challenging obstacle courses, multiple virtual environments, and checkpoints to create a structured learning experience. Users can track their performance through metrics such as completion time, checkpoint progress, and performance ratings, with potential AI-driven feedback to help identify areas for improvement. FPVR also supports registered accounts, saved performance records, attempt history, and local and global leaderboards, allowing users to monitor their progress and compare results with others. Overall, FPVR aims to make learning FPV drone piloting more accessible, engaging, and cost-effective through an immersive virtual training experience.  

\subsection{Objectives}

\wss{State what is intended to be accomplished.  The objective will be around
  the qualities that are most important for your project.  You might have
  something like: ``build confidence in the software correctness,''
  ``demonstrate adequate usability.'' etc.  You won't list all of the qualities,
  just those that are most important.}

\wss{You should also list the objectives that are out of scope.  You don't have 
the resources to do everything, so what will you be leaving out.  For instance, 
if you are not going to verify the quality of usability, state this.  It is also 
worthwhile to justify why the objectives are left out.}

\wss{The objectives are important because they highlight that you are aware of 
limitations in your resources for verification and validation.  You can't do everything, 
so what are you going to prioritize?  As an example, if your system depends on an 
external library, you can explicitly state that you will assume that external library 
has already been verified by its implementation team.}

\subsection{Custom Documents}

\wss{Summarize the CDs (Custom Documents) that were tackled by this project.  CDs 
can include usability testing, code walkthroughs, user documentation, formal
proof, GenderMag personas, Design Thinking, etc.  CDs should have already
been approved by the course instructor as included in your problem statement.
You can use a pull request to update your CDs (in TeamComposition.csv or
Repos.csv) if your plan changes as a result of the VnV planning exercise.}

\subsection{Relevant Documentation}

\wss{Reference relevant documentation.  This will definitely include your SRS
  and your other project documents (design documents, like MG, MIS, etc).  You
  can include these even before they are written, since by the time the project
  is done, they will be written.  You can create BibTeX entries for your
  documents and within those entries include a hyperlink to the documents.}

\citet{SRS}

\wss{Don't just list the other documents.  You should explain why they are relevant and 
how they relate to your VnV efforts.}

\section{Plan}

\wss{Introduce this section.  You can provide a roadmap of the sections to
  come.}

\subsection{Verification and Validation Team}

\wss{Your teammates.  Maybe your supervisor.
  You should do more than list names.  You should say what each person's role is
  for the project's verification.  A table is a good way to summarize this information.}

\subsection{SRS Verification}

\wss{List any approaches you intend to use for SRS verification.  This may
  include ad hoc feedback from reviewers, like your classmates (like your
  primary reviewer), or you may plan for something more rigorous/systematic.}

\wss{If you have a supervisor for the project, you shouldn't just say they will
read over the SRS.  You should explain your structured approach to the review.
Will you have a meeting?  What will you present?  What questions will you ask?
Will you give them instructions for a task-based inspection?  Will you use your
issue tracker?}

\wss{Maybe create an SRS checklist?}

\subsection{Design Verification}

\wss{Plans for design verification}

\wss{The review will include reviews by your classmates}

\wss{Create a checklists?}

\subsection{Verification and Validation Plan Verification}

\wss{The verification and validation plan is an artifact that should also be
verified.  Techniques for this include review and mutation testing.}

\wss{The review will include reviews by your classmates}

\wss{Create a checklists?}

\subsection{Implementation Verification}

\wss{You should at least point to the tests listed in this document and the unit
  testing plan.}

\wss{In this section you would also give any details of any plans for static
  verification of the implementation.  Potential techniques include code
  walkthroughs, code inspection, static analyzers, etc.}

\wss{The final class presentation in CAS 741 could be used as a code
walkthrough.  There is also a possibility of using the final presentation (in
CAS741) for a partial usability survey.}

\subsection{Automated Testing and Verification Tools}

\wss{What tools are you using for automated testing.  Likely a unit testing
  framework and maybe a profiling tool, like ValGrind.  Other possible tools
  include a static analyzer, make, continuous integration tools, test coverage
  tools, etc.  Explain your plans for summarizing code coverage metrics.
  Linters are another important class of tools.  For the programming language
  you select, you should look at the available linters.  There may also be tools
  that verify that coding standards have been respected, like flake9 for
  Python.}

\wss{If you have already done this in the development plan, you can point to
that document.}

\wss{The details of this section will likely evolve as you get closer to the
  implementation.}

\subsection{Software Validation}

\wss{If there is any external data that can be used for validation, you should
  point to it here.  If there are no plans for validation, you should state that
  here.}

\wss{You might want to use review sessions with the stakeholder to check that
the requirements document captures the right requirements.  Maybe task based
inspection?}

\wss{For those capstone teams with an external supervisor, the Rev 0 demo should 
be used as an opportunity to validate the requirements.  You should plan on 
demonstrating your project to your supervisor shortly after the scheduled Rev 0 demo.  
The feedback from your supervisor will be very useful for improving your project.}

\wss{For teams without an external supervisor, user testing can serve the same purpose 
as a Rev 0 demo for the supervisor.}

\wss{This section might reference back to the SRS verification section.}

\section{System Tests}

% \wss{There should be text between all headings, even if it is just a roadmap of
% the contents of the subsections.}

System tests are tests that are used to verify that FPVR satisfies the functional
and non-functional requirements as specified in the SRS. This section will cover all
system tests that will be performed end-to-end on the application.

\subsection{Tests for Functional Requirements}

% \wss{Subsets of the tests may be in related, so this section is divided into
%   different areas.  If there are no identifiable subsets for the tests, this
%   level of document structure can be removed.}

% \wss{Include a blurb here to explain why the subsections below
%   cover the requirements.  References to the SRS would be good here.}

The following is a list of tests used to cover the functional requirements (S.2) of the SRS.
Note that tests are not mutually exclusive to a single subset of requirements and can hit
multiple similarly related major functional modules. 

\subsubsection{Realistic Flight Simulation (S.2.1)}

% \wss{It would be nice to have a blurb here to explain why the subsections below
%   cover the requirements.  References to the SRS would be good here.  If a section
%   covers tests for input constraints, you should reference the data constraints
%   table in the SRS.}

\begin{enumerate}

\item{\hypertarget{test-fr-flt-01}{}\textbf{T.FR.FLT.01}\\}
\textbf{Control:} Manual
					
\textbf{Initial State:} System is running and displayed on a flat screen/monitor or VR headset.
					
\textbf{Input:} Move the drone using the corresponding input controller.
					
\textbf{Output:} The field of view of the drone changes in response to the input of the controller and moves in the expected direction
within the environment (FR-FLT-01).

\textbf{Test Case Derivation:} Verifies FR-FLT-01.
					
\item{\hypertarget{test-fr-flt-02}{}\textbf{T.FR.FLT.02}\\}
\textbf{Control:} Automatic
					
\textbf{Initial State:} Simulated drone is at rest within an environment.
					
\textbf{Input:} Simple flight path that exercises each of the drone's control axes one by one.
					
\textbf{Output:} Drone behaves as expected and follows the path within a reasonable margin of error ($\leq 5\%$ of expected path).

\textbf{Test Case Derivation:} Verifies FR-FLT-02, FR-FLT-03, FR-FLT-05.

\textbf{How test will be performed:} Starting at rest (FR-FLT-05), the drone begins by moving up and down, 
tilting left, right, forward, backward, and rotating in place (FR-FLT-02), followed with sustained flight for a period of time (FR-FLT-03).

\item{\hypertarget{test-fr-flt-03}{}\textbf{T.FR.FLT.03}\\}
\textbf{Control:} Automatic
					
\textbf{Initial State:} Drone is suspended in flight within an environment.
					
\textbf{Input:} Drone is moving and projected to hit a collision object.
					
\textbf{Output:} Drone comes into contact with object, stops moving (FR-FLT-04), and 
triggers exactly one logged crash event (FR-FLT-06).

\textbf{Test Case Derivation:} Verifies FR-FLT-04, FR-FLT-06.

\end{enumerate}

\subsubsection{Menu Interface (S.2.2, S.2.5)}

\begin{enumerate}
\item{\hypertarget{test-fr-mi-01}{}\textbf{T.FR.MI.01}\\}
\textbf{Control:} Manual
					
\textbf{Initial State:} UI is running and main menu is displayed.
					
\textbf{Input:} Free-flight mode is selected from the menu.
					
\textbf{Output:} Free-flight mode is loaded and drone is placed in an area with no courses, 
obstacles littered throughout the field that the drone can collide with (FR-LVL-02), and no attempt or 
timer limitations (FR-LVL-01).

\textbf{Test Case Derivation:} Verifies FR-LVL-01, FR-LVL-02.

\item{\hypertarget{test-fr-mi-02}{}\textbf{T.FR.MI.02}\\}
\textbf{Control:} Manual
					
\textbf{Initial State:} UI is running and main menu is displayed.
					
\textbf{Input:} User scrolls through the list of available courses, terrains/environments for selected course, 
and hits the start button.
					
\textbf{Output:} Course is loaded and drone is initialized at the starting position of the course with all stale 
data cleared (FR-LVL-07).

\textbf{Test Case Derivation:} Verifies FR-LVL-07. Navigating through the course menu also indirectly verifies FR-LVL-03,
FR-LVL-04, FR-LVL-05, FR-LVL-06, FR-UI-01, FR-UI-05.

\item{\hypertarget{test-fr-mi-03}{}\textbf{T.FR.MI.03}\\}
\textbf{Control:} Manual
					
\textbf{Initial State:} System is running and drone is in a course.
					
\textbf{Input:} Button is pressed on the controller that brings up the help menu/controller guide.
					
\textbf{Output:} Small window pops up on the screen with a list of controls and their corresponding actions (FR-UI-04). 

\textbf{Test Case Derivation:} Verifies FR-UI-04.
\end{enumerate}

\subsubsection{Checkpoint and Crash Recovery (S.2.3)}

\begin{enumerate}
\item{\hypertarget{test-fr-chk-01}{}\textbf{T.FR.CHK.01}\\}
\textbf{Control:} Automatic

\textbf{Initial State:} A selectable obstacle-course level is loaded with its
course definition available for inspection.

\textbf{Input:} Pass through checkpoints in order, then attempt to pass through
a later checkpoint before the next required checkpoint and finally reach the
finish.

\textbf{Output:} Level will record the number of checkpoints passed up until a checkpoint is skipped
from which it will stop counting (FR-CHK-02, FR-CHK-03).

\textbf{Test Case Derivation:} Verifies FR-CHK-01, FR-CHK-02, FR-CHK-03, FR-CHK-08.

\item{\hypertarget{test-fr-chk-02}{}\textbf{T.FR.CHK.02}\\}
\textbf{Control:} Automatic

\textbf{Initial State:} A course attempt has completed zero checkpoints, and a
second attempt has completed at least one checkpoint.

\textbf{Input:} Drone crashes, select respawn. Drone crashes again, select restart level.

\textbf{Output:} Application displays option to respawn and 
restart level if checkpoints $\geq 1$ (FR-CHK-04).

\textbf{Test Case Derivation:} Verifies FR-CHK-04, FR-CHK-05, FR-CHK-06, FR-CHK-07.
\end{enumerate}

\subsubsection{Timing and Performance Evaluation (S.2.4)}

\begin{enumerate}
\item{\hypertarget{test-fr-per-01}{}\textbf{T.FR.PER.01}\\}
\textbf{Control:} Automatic

\textbf{Initial State:} A new obstacle-course attempt is ready to start.

\textbf{Input:} Run through multiple checkpoint segments, crash during one
segment, record epoch time to completion.

\textbf{Output:} Total timer measures active flight, retains time from
crashed segments, stops during the recovery menu, does not reset on Respawn,
and stops at completion. The section timer measures the current segment and
resets at checkpoint completion and Respawn. The crash count increases once per
crash (FR-PER-01--FR-PER-07, FR-PER-10).

\textbf{Test Case Derivation:} Verifies FR-PER-01, FR-PER-02, FR-PER-03, FR-PER-04, FR-PER-05, FR-PER-06, FR-PER-07, FR-PER-10.

\textbf{How test will be performed:} Check $\sum_{i=1}^{n} \epsilon_{s,i} \approx \epsilon_t$ where $\epsilon_{s,i}$ is the section time 
for segment $i$ and $\epsilon_t$ is the total time.

\item{\hypertarget{test-fr-per-02}{}\textbf{T.FR.PER.02}\\}
\textbf{Control:} Automatic

\textbf{Initial State:} Drone has finished a course attempt.

\textbf{Input:} Complete attempts with known values of $T$ and $C$, including
the score boundaries 0, 399, 400, 599, 600, 799, 800, and 1000.

\textbf{Output:} The correct score and rating are calculated and displayed upon completion of a course (FR-PER-08, FR-PER-09).

\textbf{How test will be performed:} The score equals
$\max(0,1000-\lfloor T\rfloor-100C)$ and the rating is Needs Improvement for
0--399, Bronze for 400--599, Silver for 600--799, and Gold for 800 or more.

\textbf{Test Case Derivation:} Verifies FR-PER-08, FR-PER-09.
\end{enumerate}

\subsubsection{Player Accounts and Authentication (S.2.6)}

\begin{enumerate}
\item{\hypertarget{test-fr-acc-01}{}\textbf{T.FR.ACC.01}\\}
\textbf{Control:} Manual

\textbf{Initial State:} The application has network connectivity and the
username is not registered.

\textbf{Input:} Register a new username with a valid password, repeat with the
same username, and repeat with passwords shorter than 12 characters.

\textbf{Output:} The valid registration succeeds. Duplicate usernames and
short passwords are rejected with explanatory messages (FR-ACC-01, FR-ACC-02).

\textbf{Test Case Derivation:} Verifies FR-ACC-01, FR-ACC-02.

\item{\hypertarget{test-fr-acc-02}{}\textbf{T.FR.ACC.02}\\}
\textbf{Control:} Manual

\textbf{Initial State:} A registered account exists.

\textbf{Input:} Sign in with valid credentials, sign out, attempt sign-in with
an invalid username and with an invalid password, and try to submit a result
while signed out.

\textbf{Output:} Valid credentials create a session, sign-out ends it, invalid
credentials are rejected without identifying which field was wrong, and
unauthenticated practice remains available while global submission requires
authentication (FR-ACC-03--FR-ACC-05).

\textbf{Test Case Derivation:} Verifies FR-ACC-03, FR-ACC-04, FR-ACC-05.
\end{enumerate}

\subsubsection{Result Storage and Leaderboards (S.2.7)}

\begin{enumerate}
\item{\hypertarget{test-fr-dat-01}{}\textbf{T.FR.DAT.01}\\}
\textbf{Control:} Automatic

\textbf{Initial State:} A completed course result is available on an
installation with no prior result for that level.

\textbf{Input:} Save results with different scores, times, crash counts, and
timestamps, then open the local leaderboard.

\textbf{Output:} The local record contains the level, timestamp, total time,
crash count, score, rating, and local player identifier. The leaderboard shows
at least the top ten available results and orders them by descending score,
shorter time, then earlier timestamp; local-only and synchronized states are
distinguished (FR-DAT-01, FR-DAT-04, FR-DAT-06, FR-DAT-07, FR-DAT-10).

\textbf{Test Case Derivation:} Verifies FR-DAT-01, FR-DAT-04, FR-DAT-06, FR-DAT-07, FR-DAT-10.

\item{\hypertarget{test-fr-dat-02}{}\textbf{T.FR.DAT.02}\\}
\textbf{Control:} Automatic

\textbf{Initial State:} An authenticated player has completed a course.

\textbf{Input:} Submit once with connectivity, repeat while the online service
is unavailable, restore connectivity, and retry synchronization. Open both
leaderboards while the service is unavailable.

\textbf{Output:} A connected result is associated with the player's account and
appears globally. A failed submission remains local and can be synchronized
after connectivity returns. Service unavailability is reported while the local
leaderboard and training modes remain usable (FR-DAT-02, FR-DAT-03,
FR-DAT-05, FR-DAT-08).

\textbf{Test Case Derivation:} Verifies FR-DAT-02, FR-DAT-03, FR-DAT-05, FR-DAT-08.

\item{\hypertarget{test-fr-dat-03}{}\textbf{T.FR.DAT.03}\\}
\textbf{Control:} Manual

\textbf{Initial State:} A result has been synchronized for a registered player.

\textbf{Input:} Sign out, sign in on another supported installation, and open
the player's stored results.

\textbf{Output:} The previously synchronized result is retrievable and shows
the level, timestamp, total time, crash count, score, rating, and player
identifier (FR-DAT-09, FR-REV-05).

\textbf{Test Case Derivation:} Verifies FR-DAT-09, FR-REV-05.
\end{enumerate}

\subsubsection{Post-Flight Review and Intelligent Feedback (S.2.8, S.2.9)}

\begin{enumerate}
\item{\hypertarget{test-fr-rev-01}{}\textbf{T.FR.REV.01}\\}
\textbf{Control:} Manual

\textbf{Initial State:} The user has completed an obstacle-course level.

\textbf{Input:} Observe the review, retry the same level, choose another level,
and return to the main menu.

\textbf{Output:} The review appears immediately and shows the level, final
completion time, completed checkpoints, crash count, score, and rating. Each
navigation choice works without restarting the application (FR-REV-01--
FR-REV-03).

\textbf{Test Case Derivation:} Verifies FR-REV-01, FR-REV-02, FR-REV-03.

\item{\hypertarget{test-fr-rev-02}{}\textbf{T.FR.REV.02}\\}
\textbf{Control:} Automatic

\textbf{Initial State:} A course can be completed while the network and
intelligent feedback service are independently available or unavailable.

\textbf{Input:} Complete attempts in each service state and inspect the
post-flight review.

\textbf{Output:} The standard review is always displayed. When intelligent
feedback is available, it is shown automatically and its natural-language
observation is grounded in the attempt's time, checkpoint, and crash data.
Feedback failure or unavailability does not block the standard review
(FR-REV-04, FR-AI-01--FR-AI-03).

\textbf{Test Case Derivation:} Verifies FR-REV-04, FR-AI-01, FR-AI-02, FR-AI-03.
\end{enumerate}

\subsection{Tests for Nonfunctional Requirements}

\wss{The nonfunctional requirements for accuracy will likely just reference the
  appropriate functional tests from above.  The test cases should mention
  reporting the relative error for these tests.  Not all projects will
  necessarily have nonfunctional requirements related to accuracy.}

\wss{For some nonfunctional tests, you won't be setting a target threshold for
passing the test, but rather describing the experiment you will do to measure
the quality for different inputs.  For instance, you could measure speed versus
the problem size.  The output of the test isn't pass/fail, but rather a summary
table or graph.}

\wss{Tests related to usability could include conducting a usability test and
  survey.  The survey will be in the Appendix.}

\wss{Static tests, review, inspections, and walkthroughs, will not follow the
format for the tests given below.}

\wss{If you introduce static tests in your plan, you need to provide details.
How will they be done?  In cases like code (or document) walkthroughs, who will
be involved? Be specific.}

\subsubsection{Area of Testing1}
		
\paragraph{Title for Test}

\begin{enumerate}

\item{test-id1\\}

Type: Functional, Dynamic, Manual, Static etc.
					
Initial State: 
					
Input/Condition: 
					
Output/Result: 
					
How test will be performed: 
					
\item{test-id2\\}

Type: Functional, Dynamic, Manual, Static etc.
					
Initial State: 
					
Input: 
					
Output: 
					
How test will be performed: 

\end{enumerate}

\subsubsection{Area of Testing2}

...

\subsection{Traceability Between Test Cases and Requirements}

\begin{table}[htbp]
\centering
\caption{Traceability between functional requirements and system tests}
\label{tab:fr-traceability}
\begin{tabularx}{\textwidth}{p{0.32\textwidth}X}
\toprule
\textbf{System test case ID} & \textbf{Requirements verified}\\
\midrule
\hyperlink{test-fr-flt-01}{T.FR.FLT.01} & FR-FLT-01\\
\hyperlink{test-fr-flt-02}{T.FR.FLT.02} & FR-FLT-02, FR-FLT-03, FR-FLT-05\\
\hyperlink{test-fr-flt-03}{T.FR.FLT.03} & FR-FLT-04, FR-FLT-06\\
\hyperlink{test-fr-mi-01}{T.FR.MI.01} & FR-LVL-01, FR-LVL-02\\
\hyperlink{test-fr-mi-02}{T.FR.MI.02} & FR-LVL-03--07, FR-UI-01, FR-UI-05\\
\hyperlink{test-fr-mi-03}{T.FR.MI.03} & FR-UI-04\\
\hyperlink{test-fr-chk-01}{T.FR.CHK.01} & FR-CHK-01--03, FR-CHK-08\\
\hyperlink{test-fr-chk-02}{T.FR.CHK.02} & FR-CHK-04--07\\
\hyperlink{test-fr-per-01}{T.FR.PER.01} & FR-PER-01--07, FR-PER-10\\
\hyperlink{test-fr-per-02}{T.FR.PER.02} & FR-PER-08, FR-PER-09\\
\hyperlink{test-fr-acc-01}{T.FR.ACC.01} & FR-ACC-01, FR-ACC-02\\
\hyperlink{test-fr-acc-02}{T.FR.ACC.02} & FR-ACC-03--05\\
\hyperlink{test-fr-dat-01}{T.FR.DAT.01} & FR-DAT-01, FR-DAT-04, FR-DAT-06--07, FR-DAT-10\\
\hyperlink{test-fr-dat-02}{T.FR.DAT.02} & FR-DAT-02--03, FR-DAT-05, FR-DAT-08\\
\hyperlink{test-fr-dat-03}{T.FR.DAT.03} & FR-DAT-09, FR-REV-05\\
\hyperlink{test-fr-rev-01}{T.FR.REV.01} & FR-REV-01--03\\
\hyperlink{test-fr-rev-02}{T.FR.REV.02} & FR-REV-04, FR-AI-01--03\\
\bottomrule
\end{tabularx}
\end{table}

\section{Unit Test Description}

\wss{This section should not be filled in until after the MIS (detailed design
  document) has been completed.}

\wss{Reference your MIS (detailed design document) and explain your overall
philosophy for test case selection.}  

\wss{To save space and time, it may be an option to provide less detail in this section.  
For the unit tests you can potentially layout your testing strategy here.  That is, you 
can explain how tests will be selected for each module.  For instance, your test building 
approach could be test cases for each access program, including one test for normal behaviour 
and as many tests as needed for edge cases.  Rather than create the details of the input 
and output here, you could point to the unit testing code.  For this to work, you code 
needs to be well-documented, with meaningful names for all of the tests.}

\subsection{Unit Testing Scope}

\wss{What modules are outside of the scope.  If there are modules that are
  developed by someone else, then you would say here if you aren't planning on
  verifying them.  There may also be modules that are part of your software, but
  have a lower priority for verification than others.  If this is the case,
  explain your rationale for the ranking of module importance.}

\subsection{Tests for Functional Requirements}

\wss{Most of the verification will be through automated unit testing.  If
  appropriate specific modules can be verified by a non-testing based
  technique.  That can also be documented in this section.}

\subsubsection{Module 1}

\wss{Include a blurb here to explain why the subsections below cover the module.
  References to the MIS would be good.  You will want tests from a black box
  perspective and from a white box perspective.  Explain to the reader how the
  tests were selected.}

\begin{enumerate}

\item{test-id1\\}

Type: \wss{Functional, Dynamic, Manual, Automatic, Static etc. Most will
  be automatic}
					
Initial State: 
					
Input: 
					
Output: \wss{The expected result for the given inputs}

Test Case Derivation: \wss{Justify the expected value given in the Output field}

How test will be performed: 
					
\item{test-id2\\}

Type: \wss{Functional, Dynamic, Manual, Automatic, Static etc. Most will
  be automatic}
					
Initial State: 
					
Input: 
					
Output: \wss{The expected result for the given inputs}

Test Case Derivation: \wss{Justify the expected value given in the Output field}

How test will be performed: 

\item{...\\}
    
\end{enumerate}

\subsubsection{Module 2}

...

\subsection{Tests for Nonfunctional Requirements}

\wss{If there is a module that needs to be independently assessed for
  performance, those test cases can go here.  In some projects, planning for
  nonfunctional tests of units will not be that relevant.}

\wss{These tests may involve collecting performance data from previously
  mentioned functional tests.}

\subsubsection{Module ?}
		
\begin{enumerate}

\item{test-id1\\}

Type: \wss{Functional, Dynamic, Manual, Automatic, Static etc. Most will
  be automatic}
					
Initial State: 
					
Input/Condition: 
					
Output/Result: 
					
How test will be performed: 
					
\item{test-id2\\}

Type: Functional, Dynamic, Manual, Static etc.
					
Initial State: 
					
Input: 
					
Output: 
					
How test will be performed: 

\end{enumerate}

\subsubsection{Module ?}

...

\subsection{Traceability Between Test Cases and Modules}

\wss{Provide evidence that all modules have been considered.}
				
\bibliographystyle{plainnat}

\bibliography{../../refs/References}

\newpage

\section{Appendix}

This is where you can place additional information.

\subsection{Symbolic Parameters}

The definition of the test cases will call for SYMBOLIC\_CONSTANTS.
Their values are defined in this section for easy maintenance.

\subsection{Usability Survey Questions?}

\wss{This is a section that would be appropriate for some projects.}

\newpage{}

\section{Appendix --- Generative AI Use Disclosure}

\wss{ChatGPT (version 5) was used to brainstorm ideas for the template for this disclose section.}

\subsection{AI Tools Used}

\wss{Give the name of the AI tool(s) used and the version.}

\subsection{How and Where Used}

\wss{Explain what the tool(s) were used for.  Examples include: brainstorming,
explaining a technical concept, reviewing the document, generating or modifying
text, generating or modifying diagrams, identifying errors or omissions or
inconsistencies.}

\wss{For each of the uses, identify which parts of the document were affected.}

\wss{You do not need to report every interaction, but you should disclose the 
uses that materially contributed to your work}

\subsection{Verification of AI-Generated Content}

\wss{\begin{itemize}
    \item Describe how the team checked AI-generated or AI-assisted material
    before including it in the document.
    \item In particular, verify factual claims, technical assertions,
    requirements, calculations, references, and other information for which an
    AI system may produce incorrect or fabricated results.
    \item \textbf{The team is responsible for the accuracy, completeness,
consistency, and appropriateness of the submitted document, regardless of
whether material was generated by a human or an AI system.}
\end{itemize}}

\newpage{}

\section{Appendix --- Reflection}

\wss{This section is not required for CAS 741}

The information in this section will be used to evaluate the team members on the
graduate attribute of Lifelong Learning.

\input{../Reflection.text}

\begin{enumerate}
  \item What went well while writing this deliverable? 
  \item What pain points did you experience during this deliverable, and how
    did you resolve them?
  \item What knowledge and skills will the team collectively need to acquire to
  successfully complete the verification and validation of your project?
  Examples of possible knowledge and skills include dynamic testing knowledge,
  static testing knowledge, specific tool usage, Valgrind etc.  You should look to
  identify at least one item for each team member.
  \item For each of the knowledge areas and skills identified in the previous
  question, what are at least two approaches to acquiring the knowledge or
  mastering the skill?  Of the identified approaches, which will each team
  member pursue, and why did they make this choice?
\end{enumerate}

\end{document}